\documentclass[10pt]{article}

\usepackage[english]{babel}
\usepackage[letterpaper,top=0.7in,bottom=0.7in,left=0.75in,right=0.75in]{geometry}
\usepackage[T1]{fontenc}
\usepackage{lmodern}
\usepackage{microtype}
\usepackage{parskip}
\usepackage{enumitem}
\usepackage[colorlinks=true, allcolors=blue]{hyperref}

\setlength{\parskip}{0.45em}
\setlength{\parindent}{0pt}
\setlist[itemize]{leftmargin=1.2em, itemsep=0.2em, topsep=0.25em}
\pagestyle{empty}

\begin{document}

\begin{center}
{\LARGE \textbf{AIQ: The Human Capability to Harness AI}}\\[0.35em]
{\large Toward a clearer account of human--AI capability}
\end{center}

\textbf{Why now.}
AI systems are already strong enough to transform how people work, learn, and make decisions. Yet the value they create does not depend on model capability alone. Recent attention has focused on prompt engineering, context engineering, and harness engineering: how to instruct models, structure context, and build systems around them. We argue that an equally important variable remains underdeveloped: the \textbf{human capability} to use AI well. We call this capability \textbf{AIQ} -- Artificial Intelligence Quotient.

\textbf{What is AIQ.}
AIQ is a person's ability to understand where AI is strong, where it is weak, and how to design productive human--AI workflows around that moving boundary. In practice, high AIQ means knowing which parts of a workflow can genuinely benefit from AI, what context the model needs, where failure is likely, and where verification or human judgment must remain in the loop. In that sense, AIQ extends the familiar idea of \emph{human in the loop}: not just inserting a human at the end, but deliberately deciding where human judgment should guide task choice, context provision, review, escalation, and final accountability. AIQ is therefore not just about writing better prompts. It is about combining knowledge of one's own work with knowledge of AI's jagged, shifting capability frontier, while also accounting for the usability of the AI system itself: how well the system supports users in understanding what it needs, what it can do, and when its output should not be trusted.

\textbf{Why it matters.}
AI is powerful, but uneven. It can solve difficult tasks surprisingly well, then fail on nearby tasks that appear simple to humans. It also often produces plausible outputs without clearly signaling uncertainty. This makes misuse easy to miss. People with low AIQ may over-trust AI, underuse it, give it incomplete context, or deploy it at the wrong step of a workflow, including on tasks that are not well suited to current models. Conversely, even capable users can struggle when AI tools have poor usability: unclear affordances, weak support for prompt understanding, or brittle behavior under small changes in wording or context. People with high AIQ can instead turn AI into reliable augmentation: faster research, better drafting, stronger analysis, more thoughtful iteration, and improved quality control. As AI adoption accelerates, the difference between low and high AIQ may become a major driver of productivity, safety, and inclusion.

\textbf{Core dimensions of AIQ.}
If AIQ is to be useful as a concept, it should capture several distinct but connected capacities:
\begin{itemize}
\item \textbf{Task selection}: recognizing which tasks, and which steps within a task, are appropriate for AI assistance;
\item \textbf{Context provision}: supplying the relevant information, constraints, and objectives needed for good output;
\item \textbf{Prompt understanding}: knowing how different instructions shape model behavior, and what the model is likely to misunderstand;
\item \textbf{Failure anticipation and verification}: expecting hallucinations, brittleness, or misplaced confidence, and checking accordingly;
\item \textbf{Usability awareness}: recognizing that outcomes depend not only on the user, but also on whether the system helps the user understand its capabilities and limits;
\item \textbf{Robustness and generalization}: preferring workflows that remain useful across reasonable changes in phrasing, users, and context rather than relying on a single fragile prompt.
\end{itemize}
These dimensions make AIQ broader than prompt engineering and narrower than a vague claim about ``AI literacy.'' They frame AIQ as an actionable form of judgment about how humans and AI should work together.

\textbf{A broader view.}
AIQ should not be treated as a fixed personal trait or as a synonym for prompt engineering. It is better understood as an emergent capability arising from the interaction between human judgment, workflow design, context quality, and the usability and robustness of AI systems. A strong AIQ framework would therefore help explain both why the same model performs very differently across users and why better systems are not only more capable, but easier to use well.

\vfill
{\footnotesize
AIQ offers a way to move the discussion from ``what can the model do?'' to ``how do humans and AI create reliable capability together?''}

\end{document}
